\chapter{GlobalPlatform Object and Administration Model}

This chapter defines the object, authority, and lifecycle model against which
oXiDe SE is designed. Chapter~3 introduced the protocol vocabulary---APDUs,
GlobalPlatform commands, selection, and Secure Channels. The present chapter
answers the corresponding stateful questions: \emph{which resident entities
are addressed by those commands, how are they related, and under whose
authority may they be managed?}

Three kinds of statement are kept distinct throughout the chapter:

\begin{description}
\item[GlobalPlatform reference model.] These paragraphs summarize the
      concepts and command semantics defined by the GlobalPlatform Card
      Specification.
\item[oXiDe SE implementation contract.] These paragraphs state the
      architecture and invariants that an oXiDe SE implementation is expected
      to enforce. They are design requirements, not claims of formal proof.
\item[Beta implementation status.] These paragraphs identify the implemented
      subset and the remaining validation or implementation work. Detailed
      mechanisms and evidence are deferred to the implementation and test
      chapters.
\end{description}

This separation is important. Compatibility with the GlobalPlatform object
model does not require oXiDe SE to reproduce every optional management mode,
and an internal oXiDe SE representation must not be mistaken for a structure
mandated by GlobalPlatform.

%--------------------------------------------------
\section{Platform Object Model}
\label{sec:gp-object-model}

% Plan:
% - define the main GP-managed entities
% - explain that AIDs are the stable identifiers used by registry and dispatch
% - distinguish executable content, installed instances, and security authorities
% - introduce the registry view as the structured catalog of those entities

\subsection{GlobalPlatform Reference Entities}

GlobalPlatform does not begin from a process model or an operating-system
abstraction. It begins from managed entities resident on the secure element
and from the relationships between them. The principal reference entities are
\gpref{gp-card}{§7, §9, §11}:

\begin{itemize}
\item an \textbf{Executable Load File}, which represents deployable code
  transported to the secure element through the loading protocol;
\item an \textbf{Executable Module}, which identifies installable code within
  an Executable Load File;
\item an \textbf{Application}, or more precisely an installed application
  instance, which may later be selected and addressed through APDUs;
\item a \textbf{Security Domain}, which is also an application-like
  resident entity but with management responsibilities, privileges, and
  associated cryptographic material;
\item the unique \textbf{Issuer Security Domain} of one GlobalPlatform
  environment, which is the first application and has the special role
  defined by the Card Specification.
\end{itemize}

These entities are not merely names in a registry. They are the objects
whose lifecycle is controlled by commands such as \texttt{LOAD},
\texttt{INSTALL}, \texttt{DELETE}, \texttt{SET STATUS}, and
\texttt{GET STATUS}. The command vocabulary introduced in Chapter~3
therefore only becomes meaningful once these entities have been
distinguished.

Conceptually, the \textbf{GlobalPlatform Registry} is the organized view of
the managed world inside the secure element. It records which load files,
modules, applications, and Security Domains exist, how they are associated,
and which privileges and lifecycle states apply. Security-Domain keys and data
belong to the administrative model, but representing each key as an
AID-addressable registry object is not required by GlobalPlatform.

\subsection{oXiDe SE Managed Objects}

oXiDe SE implements one persistent registry containing packages, Rustlet
instances, Security Domain instances, key objects, and generic data objects.
Using the same bounded representation for these families is an oXiDe SE design
choice. It permits the same ownership, persistence, lifecycle, and deletion
rules to be applied consistently, while command handlers retain the distinct
GlobalPlatform semantics of each object family.

One oXiDe package record represents one FAE Executable Load File and records
the single Executable Module AID from which its Rustlet instances may be
created. Executable Modules are therefore projected in \texttt{GET STATUS},
but are not independently stored or lifecycle-managed registry entries. A
multi-module Executable Load File is outside the beta model.

Every managed object is attached to an administrative Security Domain. The
internal field \path{parent_sd_aid} records that containment relation. One
exception is structural: the oXiDe Root Security Domain has no parent and is
the unique root of the internal registry tree. This root is an oXiDe SE
platform entity; it is not necessarily the GlobalPlatform Issuer Security
Domain.

%--------------------------------------------------
\subsection{Object Identification}
\label{sec:gp-object-identification}

% Plan:
% - explain AIDs as the primary GP identifiers
% - distinguish load-file AID, application-instance AID, and security-domain AID
% - state uniqueness and role in selection, management, and registry references

The primary identifier used throughout GlobalPlatform is the
\textbf{Application Identifier} (AID). Even though the name comes from
ISO/IEC~7816 and from the application-selection model, GlobalPlatform
uses AIDs more broadly: an AID identifies not only a selectable
application instance, but also Security Domains and load-file-related
objects referenced by management commands \gpref{gp-card}{§7, §11}.

Several distinct AIDs may therefore coexist for conceptually related
objects:

\begin{itemize}
\item an \textbf{Executable Load File AID} identifies the executable content being
  transported and installed;
\item an \textbf{Application AID} identifies an installed instance that
  may later become selectable;
\item a \textbf{Security Domain AID} identifies the administrative
  entity under whose authority management actions are performed.
\end{itemize}

This distinction matters because GlobalPlatform deliberately separates
deployable code, installed instances, and administrative authority. A
single load file may give rise to one or more installed instances, and
those instances may later be placed under the responsibility of a given
Security Domain. The identifiers must therefore remain stable across
management operations and allow those relationships to be expressed
explicitly.

From a protocol perspective, AIDs are also the bridge between the APDU
world and the object model. They appear in \texttt{SELECT}, in
\texttt{INSTALL} parameters, in registry-oriented queries, and in
administrative operations such as deletion or lifecycle updates. An AID
is therefore not just a name looked up at selection time; it is the
stable handle by which GlobalPlatform refers to a managed secure-element
entity across
its whole management model.

For an off-card operator, a minimal useful set of identifiers is
therefore:
\begin{itemize}
\item the AID of the oXiDe Root Security Domain when a platform-wide
  administration operation is intentionally requested;
\item the AID of the Issuer Security Domain, when one is present, which is the
  usual entry point for secure-element administration;
\item the AID of any supplementary Security Domain used for delegated or
  specialized administration;
\item the AID of the target application instance to be selected or
  personalized;
\item the AID of the corresponding load file when deployment or removal
  operations are performed.
\end{itemize}

The Root Security Domain AID and an Issuer Security Domain AID designate
different roles even when a particular configuration chooses the same backend
implementation for both. The former selects the oXiDe SE platform root; the
latter, when declared, selects the unique Issuer Security Domain of the
GlobalPlatform environment.

%--------------------------------------------------
\subsection{Object Relationships}
\label{sec:gp-object-relationships}

% Plan:
% - explain dependency graph between load files, apps, and security domains
% - explain associated security-domain relationship
% - introduce the idea that management invariants arise from these relations

The GlobalPlatform object model is relational rather than flat. The most
important relationships are:

\begin{itemize}
\item an installed application instance depends on executable content
  that was previously made present on the secure element through the loading
  process;
\item an application instance is associated with a Security Domain that
  acts as its administrative authority;
\item a Security Domain owns keys, privileges, and management authority
  that are distinct from the application business logic itself.
\end{itemize}

These relationships are not optional conceptual decorations. They are
what gives management commands their preconditions and their effects. An
application cannot be installed if the relevant executable content is not
present; a simple load-file deletion cannot proceed while installed instances
still depend on it, whereas cumulative deletion explicitly includes related
objects; a sensitive management action is interpreted under the authority
of a Security Domain rather than under the identity of an arbitrary
selected application.

\subsection{GlobalPlatform Association Forest}

GlobalPlatform associates each Application and Executable Load File with a
Security Domain. Because a Security Domain is itself an Application, Security
Domains may form association hierarchies. The Issuer Security Domain is
effectively associated with itself. GlobalPlatform also permits another
Security Domain to become self-associated through extradition, thereby forming
the root of another association hierarchy \gpref{gp-card}{§7.1, §7.2}.

There is no contradiction between one Issuer Security Domain and several
association hierarchies. Only one root has the special Issuer Security Domain
role. Other self-associated roots remain supplementary Security Domains. Their
providers may use distinct keys and administer content independently of the
Card Issuer's ordinary management context.

The GlobalPlatform association model is therefore a forest. Its usual shape
has the Issuer Security Domain at the root of the principal hierarchy, with
most supplementary Security Domains below it. Additional self-associated
Security Domains may head detached hierarchies.

\subsection{oXiDe SE Internal Root and Tree Projection}

oXiDe SE represents that forest as one internal tree by adding a unique
technical root:

\begin{verbatim}
oXiDe Root Security Domain
|-- Issuer Security Domain
|   |-- ordinary Security Domain A
|   `-- ordinary Security Domain B
`-- detached Security Domain C
    `-- ordinary Security Domain C.1
\end{verbatim}

In a production-oriented composition, at least the oXiDe Root Security Domain
and the Issuer Security Domain are normally pre-deployed. The Issuer Security
Domain is explicitly marked as such and is a direct child of the root;
ordinary Security Domains are normally attached to the Issuer Security Domain
or one of its descendants. A GlobalPlatform hierarchy detached from the Issuer
hierarchy is represented internally by attaching its self-associated root
directly to the oXiDe root. Development compositions may deliberately declare
only the technical root and therefore have no distinct Issuer Security Domain.

The first internal edge has a different interpretation from an ordinary
GlobalPlatform association. In the GlobalPlatform-facing projection, the
Issuer Security Domain and every detached hierarchy root are self-associated.
Their internal parent is the oXiDe root only for platform containment,
visibility, persistence, and recovery. Registry responses must preserve this
distinction and must not accidentally expose the technical root as the
GlobalPlatform associated Security Domain of those roots.

\subsection{Uniform Subtree Visibility}

oXiDe SE gives the internal tree one uniform operational interpretation. A
Security Domain may resolve its own objects and every object contained in its
descendant Security Domains. It cannot resolve objects owned by an ancestor or
by a sibling branch. Objects outside that view are treated as absent for the
administrative operation, which avoids disclosing their existence.

For every Security Domain $S$, let $V(S)$ be its visible registry set. If $C$
is a child of $P$, the structural invariant is:

\begin{center}
$V(C) \subseteq V(P)$.
\end{center}

The root's descendants are all managed objects, so $V(root)$ is the complete
registry. The Issuer Security Domain sees the usual Issuer hierarchy but not a
detached root or its descendants. Two detached hierarchies are mutually
invisible. A newly created object is attached to the active Security Domain
that authorizes its creation. In the beta, detached hierarchy roots may be
declared by predeployment beneath the technical root; post-issuance
extradition and reattachment are not implemented.

Visibility is necessary but not sufficient for authorization. Command type,
delegated privileges, Secure Channel state, lifecycle, dependency checks, and
the selected Security Domain's own policy remain mandatory after lookup.

\subsection{Uniform Authority and Monotone Delegation}
\label{sec:gp-privilege-delegation}

When this chapter speaks about Security-Domain ``rights'', it does not mean
ordinary application permissions. It means GlobalPlatform administrative
privileges: the authority to perform families of management operations on
resident objects. Examples include loading executable content, installing an
application instance, deleting managed objects, changing lifecycle state,
managing secure-channel keys, performing delegated management, requesting or
verifying deployment tokens, generating receipts, or exercising card-level
lock and termination operations.

These privileges are attached to Security Domain instances. An ordinary
application instance does not become a management authority merely because it
lives under a privileged Security Domain. It is administered by that Security
Domain, and it may rely on services made available through that association,
but registry mutation and platform administration remain Security-Domain
authority.

GlobalPlatform defines specialized privileges whose roles are not, by
themselves, a simple parent-to-child lattice. oXiDe SE deliberately imposes a
stricter and more uniform rule: a Security Domain may delegate only authority
that it already possesses, and a child must not receive a broader privilege
set than its parent. If $R(S)$ denotes the operations granted to $S$, then:

\begin{center}
$R(C) \subseteq R(P)$.
\end{center}

Combined with subtree visibility, this gives
$V(C) \subseteq V(P)$ and $R(C) \subseteq R(P)$. Administrative delegation
may restrict or subdivide authority; it cannot amplify either its operations
or its reach.

GlobalPlatform's \emph{Global Delete}, \emph{Global Lock}, and
\emph{Global Registry} are privileges, not distinct instruction bytes. In the
reference model, an Application or Security Domain holding the corresponding
privilege may act independently of ordinary Security-Domain association or
hierarchy restrictions \gpref{gp-card}{§9.1.3.4--§9.1.3.5, §9.6.5}.

oXiDe SE intentionally adopts a narrower beta contract. Only the technical
root obtains transverse reach, as a consequence of being the ancestor of every
internal hierarchy. A non-root Security Domain never escapes its subtree, even
if its stored privilege bytes contain a Global Delete, Global Lock, or Global
Registry bit. The common mutation and query primitives remain parameterized by
an authority AID; invoking them with the root AID gives them complete reach
without bypassing registry, lifecycle, persistence, or dependency checks. This
is an explicit oXiDe SE profile restriction relative to the GlobalPlatform
privilege semantics, not a restatement of the standard.

%--------------------------------------------------
\subsection{Selected Context and Visible Endpoints}
\label{sec:gp-selected-context}

% Plan:
% - explain what a host actually addresses when talking to the card
% - distinguish selected Security Domain and selected ordinary application
% - connect this to Chapter 3 without entering kernel design

From the point of view of an external management tool, the secure
element does not expose a single flat command receiver. It exposes a set
of \textbf{selectable contexts}. Through \texttt{SELECT}, the off-card
entity chooses which AID becomes the current communication endpoint. Four
situations matter most:

\begin{itemize}
\item the \textbf{oXiDe Root Security Domain} is selected for an explicitly
  platform-wide operation. Its ordinary subtree is the complete registry;
\item a \textbf{Security Domain} is selected, typically the Issuer
  Security Domain, and subsequent commands are interpreted in the
  administrative plane of the secure element;
\item an \textbf{ordinary application} is selected, and subsequent
  commands are interpreted according to that application's own semantics;
\item no application-specific management action is possible until the
  appropriate administrative context has first been selected.
\end{itemize}

This explains an important operational distinction already hinted at in
Chapter~3. Commands such as \texttt{LOAD}, \texttt{INSTALL},
\texttt{DELETE}, \texttt{PUT KEY}, \texttt{GET STATUS}, or
\texttt{INITIALIZE UPDATE} are not merely generic APDU opcodes that can
be sent to any selected endpoint. They are meaningful only in the
appropriate administrative context, i.e. under a Security Domain
authorized to interpret them.

In oXiDe SE this authorization is not modeled as a collection of
independent ad hoc checks spread across the kernel. The object registry is
also the visibility boundary for the resolved administrative authority. A
lookup under one Security Domain resolves only within the subtree
administratively reachable from that Security Domain. An object outside this
visible subtree is treated as absent in that context. This design makes the
management rule above mechanically enforceable: the resolved authority may
act on its own objects and on objects of its descendants, while new objects
are created under that authority.

The beta resolves that authority differently for its two supported development
paths. A protected management APDU is bound to the active Security Domain and
its Secure Channel. A clear management APDU is bound to the technical root and
is accepted only by a permissive root backend such as
\texttt{NullSecurityDomain}. Selecting a child Security Domain does not make
clear management APDUs execute under that child. Logical-channel-specific
selection and Secure Channel contexts remain future work.

Conversely, once an application instance has been selected, the same APDU
transport is used, but the semantics changes entirely: the secure element
interprets the APDU according to the application contract rather than
according to the GlobalPlatform administration model.

%--------------------------------------------------
\section{Lifecycle Model}
\label{sec:gp-lifecycle-model}

% Plan:
% - define lifecycle as state machine over GP objects
% - explain that transitions are command-driven and constrained
% - prepare the object-specific lifecycle subsections

GlobalPlatform associates lifecycle states with the principal managed
objects on the secure element. A lifecycle state is more than a diagnostic label:
it controls what operations are currently permitted on a given entity and
which command sequences remain legal \gpref{gp-card}{§9.6, §11}.

Conceptually, the lifecycle model combines three ideas:

\begin{itemize}
\item a managed object is always in a well-defined state;
\item only certain transitions are allowed from a given state;
\item commands such as \texttt{INSTALL}, \texttt{SET STATUS},
  \texttt{DELETE}, and \texttt{SELECT} are interpreted against that
  state machine rather than in isolation.
\end{itemize}

This is why administration and deployment cannot be described purely as a
list of commands. They are organized as stateful workflows. A load-file
transfer, an application installation, a personalization session, or a
deletion sequence only make sense because the addressed object moves
through allowed states and because invalid transitions are rejected.

For a security specialist using the secure element, the practical consequence is
that the current state of the addressed object is always part of the
precondition of a command. A management script therefore needs more than
an AID and an opcode: it needs a correct view of whether the target load
file is present, whether the application instance already exists, whether
it is selectable, and whether the selected Security Domain is currently
allowed to perform the requested transition.

GlobalPlatform defines object-specific lifecycle models
\gpref{gp-card}{§5.2--§5.3}. An Executable Load File has only the
\texttt{LOADED} state. An ordinary Application has \texttt{INSTALLED},
\texttt{SELECTABLE}, and \texttt{LOCKED} states, plus any application-specific
states. A supplementary Security Domain has \texttt{INSTALLED},
\texttt{SELECTABLE}, \texttt{PERSONALIZED}, and \texttt{LOCKED} states. The
Issuer Security Domain instead inherits the Card Life Cycle State. The
following subsections do not replace those definitions. They describe the
deliberately smaller and partly implementation-specific projection used by the
oXiDe SE beta registry.

Absence, deletion, and failed installation are represented by the absence of a
committed registry entry, not by an additional lifecycle value. oXiDe SE uses
one compact state pair per ordinary managed-object family:

\begin{center}
\begin{tabular}{p{0.22\textwidth}p{0.22\textwidth}p{0.42\textwidth}}
\textbf{Object family} & \textbf{States} & \textbf{Operational guard} \\
\hline
\texttt{Package} & \texttt{Loaded}, \texttt{Locked} &
  only \texttt{Loaded} packages may be instantiated; \texttt{Locked} is an
  oXiDe extension, not a GlobalPlatform Executable Load File state; \\
\texttt{Instance} & \texttt{Selectable}, \texttt{Locked} &
  only \texttt{Selectable} instances may be selected; \\
\texttt{SecurityDomain} & \texttt{Selectable}, \texttt{Locked} &
  only \texttt{Selectable} Security Domains may be selected or used to open a
  secure channel; \\
\texttt{Key} & \texttt{Active}, \texttt{Locked} &
  only \texttt{Active} keys may be used by secure-channel derivation; these are
  internal oXiDe states, not GlobalPlatform Registry lifecycle values.
\end{tabular}
\end{center}

The internal \texttt{Locked} state is conservative. It keeps the object
visible to management code, but prevents the operation that would normally
make the object effective: instantiating a package, selecting an instance,
opening a secure channel through a Security Domain, or deriving session keys
from a key object. The data model can represent all four locked forms, but the
current APDU surface does not expose all their transitions. Future profiles
may add richer state values, but those values are not part of the beta
lifecycle contract.

The oXiDe root is not an ordinary GlobalPlatform lifecycle object: it exists
from bootstrap, has no parent, and cannot be deleted or have its state changed
through the beta management APDUs. The Issuer Security Domain is also distinct
from an ordinary supplementary Security Domain because GlobalPlatform makes it
inherit the Card Life Cycle State. The beta registry nevertheless stores both
Issuer and supplementary Security Domains using the compact
\texttt{Selectable}/\texttt{Locked} pair. Card lifecycle and the complete
Issuer lifecycle projection are not implemented.

%--------------------------------------------------
\subsection{Application Lifecycle}
\label{sec:gp-lifecycle-application}

% Plan:
% - describe the conceptual lifecycle of an application instance
% - emphasize installation, selectability, locking, and administrative control
% - link conceptually to INSTALL / SET STATUS without implementation details

For an application instance, the retained lifecycle distinction is simple:
the instance is either selectable or locked. GlobalPlatform distinguishes
installation from the transition that makes an instance selectable. The beta
profile currently uses the combined \texttt{INSTALL [for install and make
selectable]} form and creates the instance directly in the
\texttt{Selectable} state when policy, package lookup, and the Rustlet
installation call all succeed. A failed installation does not create a stable
registry object.

The practical consequence is that \texttt{SELECT} is authorized by the
instance state. A \texttt{Selectable} instance may become the current
application context; a \texttt{Locked} instance remains known to management
code but cannot be selected. Deletion removes the registry entry and therefore
returns the instance to the operationally absent condition.

The normal application-instance transitions are:

\begin{itemize}
\item \textbf{absent to selectable}: caused in the beta profile by the
  successful combined install-and-make-selectable form of \texttt{INSTALL}.
\item \textbf{selectable to locked}: represented by the beta's restricted
  \texttt{SET STATUS} path or by kernel policy.
\item \textbf{locked to selectable}: represented by the same restricted
  path. The full GlobalPlatform status-type and previous-state rules are not
  yet implemented.
\item \textbf{selectable or locked to absent}: caused by \texttt{DELETE},
  provided dependency and privilege checks allow removal.
\end{itemize}

\texttt{SELECT} itself does not install an application and does not create a
new lifecycle state. It only makes an already selectable application the
current execution context. Conversely, deselection is a change of current
context, not a registry lifecycle transition.

%--------------------------------------------------
\subsection{Load File Lifecycle}
\label{sec:gp-lifecycle-loadfile}

% Plan:
% - explain load files as deployable content with their own lifecycle
% - distinguish load-file presence from application-instance usability
% - explain dependency constraints on removal

The lifecycle of executable content is different from that of application
instances. A package is first transported, validated and committed as executable
content. Its conceptual role is to make code available on the secure element so
that one or more application instances may later be installed from it.

This means that a load file may exist even when no currently selectable
application has yet been created from it, and conversely that an
installed application may continue to depend on that load file for as
long as the instance exists. The load-file lifecycle therefore captures
availability and dependency rather than direct selectability.

From a management perspective, this is what explains the asymmetry
between \texttt{LOAD} and \texttt{DELETE}. Loading introduces executable
content as a managed dependency; deletion must later respect the fact
that other objects may still rely on that content.

The oXiDe package transitions are:

\begin{itemize}
\item \textbf{absent to loaded}: caused by a complete and valid
  \texttt{INSTALL [for load]} plus \texttt{LOAD} sequence. The temporary load
  context is volatile until the final block has been checked; only the final
  success creates the package registry entry.
\item \textbf{loaded to locked}: caused by the beta's diagnostic
  \texttt{SET STATUS} extension or kernel policy. A locked package remains
  present but cannot be used by \texttt{INSTALL [for install]} to create a new
  instance. GlobalPlatform itself defines no \texttt{LOCKED} Executable Load
  File lifecycle state.
\item \textbf{locked to loaded}: caused by the corresponding oXiDe diagnostic
  unlock operation.
\item \textbf{loaded or locked to absent}: caused by a simple \texttt{DELETE}
  when no installed instance or dependent object still requires the load file,
  or by cumulative deletion that explicitly includes the related objects.
\end{itemize}

This model deliberately keeps partially received code outside the stable
registry. The registry must never advertise a load file as usable before the
whole binary has been received and validated.

%--------------------------------------------------
\subsection{Security Domain Lifecycle}
\label{sec:gp-lifecycle-domain}

% Plan:
% - explain that security domains also have lifecycle states
% - stress their administrative role and privileged status
% - distinguish them from ordinary applications without entering implementation

Security Domains also have lifecycle states, but their meaning is more
administrative than purely functional. A Security Domain is not only a
selectable entity; it is also a carrier of privileges, keys, and
management authority. As a result, its lifecycle constrains not only its
own reachability but also which administrative operations may lawfully be
performed under its authority.

Three cases must be distinguished in oXiDe SE. The Root Security Domain is the
non-deletable technical anchor and is outside ordinary Security Domain
lifecycle transitions. When present, the Issuer Security Domain is the unique
GlobalPlatform issuer endpoint and inherits the card lifecycle in the
GlobalPlatform reference model. Supplementary and detached Security Domains
are ordinary lifecycle-bearing administrative principals.

In many operational environments, the Issuer Security Domain is also the
administrative endpoint informally associated with the \emph{Card Manager}
role. In strict GlobalPlatform terminology, however, \emph{Card Manager} is a
broader historical label than \emph{Issuer Security Domain}. This document
therefore keeps \emph{Issuer Security Domain} as the precise preferred term for
the selected administrative endpoint with which an off-card management tool
first establishes trust and through which it performs secure-element-management
operations.

For an ordinary Security Domain, the compact beta transitions are:

\begin{itemize}
\item \textbf{absent to selectable}: caused by the successful combined
  install-and-make-selectable form of \texttt{INSTALL} when the instance is declared as a
  Security Domain. The command fixes the Security Domain instance AID, backend
  or package binding, initial administrative state, and base privilege bytes.
  Once selectable, the Security Domain may be selected as an administrative
  endpoint and may establish secure channels according to its key and profile
  policy.
\item \textbf{selectable to locked}: represented by the beta's restricted
  \texttt{SET STATUS} transition or by platform policy. A locked Security
  Domain remains a registry object, but it cannot be selected and cannot be
  used to establish a secure channel. Because secure-channel establishment is
  the entry point for protected management authority, this also prevents normal
  administrative use under that Security Domain.
\item \textbf{locked to selectable}: represented by the same restricted
  unlock path.
\item \textbf{selectable or locked to absent}: caused by simple
  \texttt{DELETE} when no descendants or key dependencies remain, or by an
  authorized cumulative deletion that includes the complete subordinate
  closure.
\end{itemize}

The key difference from ordinary application lifecycle is that Security Domain
state gates authority as well as selectability. A selectable application may
process application APDUs; a selectable Security Domain may additionally
authorize administrative operations, manage key material, and host secure
channel state. A locked Security Domain therefore does not merely become
unselectable: it also stops being a usable management authority.

The lifecycle returned through the standard lifecycle data object
\texttt{9F70} must reflect the effective state of the selected Security
Domain. In oXiDe SE terms, this is the state stored with the Security Domain
registry object or, for a Rustlet-backed Security Domain, the state reported by
the selected Rustlet Security Domain through its dedicated management
interface.

The current \texttt{SET STATUS} decoder is not yet a conformant implementation
of the GlobalPlatform state machine. GlobalPlatform uses P1 values to
distinguish the Issuer Security Domain, Applications or supplementary Security
Domains, and a Security Domain together with its associated Applications; P2
contains the requested standard lifecycle value
\gpref{gp-card}{§11.10}. The beta instead accepts an internal package category,
maps its three categories directly to package, application, and Security
Domain registry kinds, and accepts only compact lock/unlock controls. In
particular it neither implements the Issuer/Card Life Cycle State nor the
standard recursive Security-Domain-and-associated-Applications form. This path
is diagnostic functionality pending protocol realignment.

%--------------------------------------------------
\subsection{Key Object State}
\label{sec:gp-lifecycle-key}

GlobalPlatform associates key material and key-management operations with a
Security Domain, but it does not define the oXiDe key entries as ordinary
GlobalPlatform Registry objects with a common lifecycle. oXiDe nevertheless
stores typed key objects beneath their owning Security Domain. Their internal
state is orthogonal to key version, key identifier, usage, and raw material and
answers one operational question: may this key be used by secure-channel
establishment?

The current externally reachable behavior is:

\begin{itemize}
\item \textbf{absent to active}: caused by successful key provisioning, such
  as \texttt{PUT KEY}. The key object becomes available to the corresponding
  secure-channel profile.
\item \textbf{active to active}: a later \texttt{PUT KEY} may replace the
  matching owner-qualified key identity.
\end{itemize}

The registry representation also contains \texttt{Locked}, and secure-channel
lookup correctly ignores a key in that state, but no conformant management APDU
currently performs the Active/Locked transition. Generic deletion by the
synthetic internal key-object AID is an implementation facility, not a
GlobalPlatform key-removal interface. Key rotation does not require a separate
lifecycle state. Moreover, a multi-key \texttt{PUT KEY} is presently applied
and persisted entry by entry: complete key-set replacement is not atomic and a
late parsing or storage failure may leave earlier entries updated. This is a
beta limitation that must be removed before claiming complete
\texttt{PUT KEY} interoperability.

%--------------------------------------------------
\section{Deployment Workflow}
\label{sec:gp-deployment-workflow}

% Plan:
% - describe the canonical deployment flow at GP level
% - keep it conceptual, not implementation-specific
% - connect secure channel, load, install, personalization, and activation

At a conceptual level, GlobalPlatform deployment is a staged workflow
rather than a single command. A typical sequence is:

\begin{enumerate}
\item select the relevant Security Domain, typically the Issuer Security
  Domain;
\item establish the administrative authority and, when required, a
  secure channel through \texttt{INITIALIZE UPDATE} and
  \texttt{EXTERNAL AUTHENTICATE};
\item prepare the secure element for receiving executable content;
\item transfer that content through one or more \texttt{LOAD} commands;
\item create or configure the installed instance through one of the
  \texttt{INSTALL} variants;
\item optionally personalize the instance;
\item make it selectable or otherwise transition it to the intended
  operational state.
\end{enumerate}

In oXiDe SE, a Rustlet normally follows this workflow after issuance: its FAE
image is built independently, transferred to an already deployed secure
element, and installed under the authority of a Security Domain. For a
controlled initial image, oXiDe SE also supports pre-issuance deployment. The
build configuration then embeds the Rustlet and declares its package and
instance relationships in the predeployment manifest, so that they form part of
the initial registry. Predeployment changes when the content enters the system;
it does not turn the Rustlet into kernel code or bypass the managed-object,
authority, lifecycle, and isolation models described in this chapter.

\subsection{Bootstrap and the Initial Administrative Tree}

A production-oriented oXiDe SE image normally pre-deploys at least two distinct
administrative objects. The Root Security Domain establishes the technical
root, initial registry authority, and recovery context. The unique Security
Domain explicitly declared with \texttt{role = "issuer"}, when present, is a
direct child of that root and is the normal entry point for subsequent
card-content management. Other pre-deployed Security Domains are ordinarily
placed below the Issuer hierarchy. A Security Domain attached directly to the
oXiDe root is projected as the self-associated root of a detached
GlobalPlatform hierarchy.

This is one schema rather than a set of image profiles. A kernel-only image
uses only the materialized technical root and kernel APDU modules. A
Rustlet-development image declares a permissive root with no child Security
Domain. A Security-Domain test image adds whatever hierarchy it exercises. A
production-oriented hierarchy declares exactly one Issuer and the required
protected policies. Zero Issuer declarations therefore means that no distinct
Issuer exists; one identifies it; a second is rejected. The build validates
parents, cycles, duplicate AIDs, and monotone privileges, then materializes the
root, the optional Issuer, and the remaining domains in parent-before-child
order independently of TOML declaration order.

This bootstrap is part of the trusted initial image. Post-issuance management
cannot manufacture a second parentless object, replace the oXiDe root, or
change its lifecycle through the beta APDU surface. Dynamic extradition and
detachment are not implemented; detached roots are currently established only
by trusted predeployment.

This workflow explains why the same command vocabulary is reused in
different combinations. For example, \texttt{INSTALL} is not a single
one-dimensional action: its mode flags let it participate in loading,
installation, personalization, make-selectable, registry-update, or
extradition workflows. The command set is therefore best read as a
workflow language for secure-element administration.

At this stage, two operational points are worth making explicit.

First, management activity is usually embedded in an authenticated
session. In a secure deployment, an operator is not expected to send
\texttt{LOAD}, \texttt{INSTALL}, \texttt{PUT KEY}, or similar commands
as isolated cleartext APDUs. The normal pattern is to select the
administrative Security Domain, establish its configured Secure Channel---for
example SCP03 or an SCP11 variant---and only then execute the management
workflow.

Second, personalization is conceptually part of deployment, but it is not
the same thing as installation. Installation creates or configures the
instance; personalization transfers instance-specific data. This is why a
deployment sequence may remain incomplete, from the operational point of
view, even after the executable content has been loaded and the
application instance has been created.

The beta profile implements the Rustlet load/install path needed by its test
and development workflows, but not every GlobalPlatform personalization,
extradition, Token, DAP, or Receipt form described by the reference workflow.
Those omissions are implementation limits, not alternative definitions of the
GlobalPlatform operations.

\subsection{Hierarchy Removal and Cumulative Delete}

Deletion is evaluated by the same authority-and-visibility rule as creation or
update. A simple deletion removes one visible object only when doing so leaves
no live dependency. Deleting a package with installed instances, or deleting a
Security Domain that still owns objects or subordinate domains, must therefore
fail in the simple form.

Cumulative deletion changes the requested target set, not the authorization
model. GlobalPlatform~2.3.1 defines hierarchy or sub-hierarchy cumulative
deletion as an optional operation performed by a Security Domain with Global
Delete privilege \gpref{gp-card}{§9.5.3--§9.5.4}. The oXiDe beta applies a
narrower tree rule: an authority may request the cumulative form only over a
target and closure already contained in its own visible subtree; transverse
scope is obtained only by using the technical root.

The implementation first computes the complete related-object closure:
instances belonging to a package, or Security Domains and managed content in a
subordinate hierarchy. Every member must remain within the active authority's
visible subtree. The in-RAM entries are removed dependants first and the
resulting registry is published in one persistent snapshot. Flash recovery can
therefore retain the previous valid snapshot if publication is interrupted;
this does not claim a general transaction mechanism for arbitrary mutations.

The oXiDe root applies this same algorithm. Because every hierarchy is in its
subtree, a root-requested cumulative deletion can remove any detached
hierarchy and has card-wide effect. No variant of \texttt{DELETE}, including a
cumulative form, may remove the oXiDe root itself. The beta does not implement
GlobalPlatform's logically deleted Executable Load File state or the complete
cross-hierarchy reference-repair behavior required by cumulative deletion.

%--------------------------------------------------
\section{Administration Model}
\label{sec:gp-administration-model}

% Plan:
% - explain administrative authority in GP terms
% - define the role of security domains and privileges
% - distinguish management APDUs from ordinary application semantics

GlobalPlatform administration is not defined as unrestricted access to a
flat command set. It is defined through \textbf{authority} and
\textbf{privileges}. Administrative actions are performed under the
responsibility of a Security Domain, and the specification associates
privileges with such entities to determine which operations they may
perform \gpref{gp-card}{§7, §9.6, §11}.

This explains why commands such as \texttt{PUT KEY}, \texttt{LOAD},
\texttt{INSTALL}, \texttt{DELETE}, \texttt{SET STATUS}, or
\texttt{GET STATUS} are not conceptually ordinary business commands.
They belong to the administrative plane of the secure element. They operate on
managed objects, on registry state, on keys, and on lifecycle
transitions. Even when they are addressed in the context of a selected
Security Domain, their semantics remain secure-element administration semantics.

The privilege model is therefore the conceptual bridge between the object
model and the command model. It explains why not every selected application
may perform every administrative action, why secure channel establishment
is tied to particular keys and domains, and why lifecycle changes are not
purely local decisions of the addressed application instance.

\subsection{Ordinary and Global Command Forms}

In the GlobalPlatform reference model, \emph{Global Delete}, \emph{Global
Lock}, and \emph{Global Registry} name privileges and their expanded effects;
they are not separate DELETE, SET STATUS, or GET STATUS instruction bytes. A
holder of the relevant global privilege may cross the association boundary
that constrains an ordinary authority. Cumulative Delete is a distinct optional
DELETE mode and is not synonymous with every use of Global Delete.

oXiDe SE does not define a separate set of mutation primitives for card-wide
administration. The internal operation is parameterized by a resolved Security
Domain authority and is evaluated against that authority's visible subtree.
At the APDU boundary the beta classifies DELETE, SET STATUS, and GET STATUS
resolved to the technical root as internal ``global'' forms, then invokes the
same deletion, status, or registry operation with \texttt{authority = root}.
There is no additional on-wire opcode. Because the root's subtree contains
every hierarchy, that invocation has card-wide effect.

A subordinate Security Domain invokes the common operation over a smaller
visible set. Stored Global Delete, Global Lock, or Global Registry bits do not
expand that set in the beta. Consequently oXiDe's root-only transverse model is
stricter than, and not fully interoperable with, the GlobalPlatform global
privilege model. Neither root scope nor the internal classification bypasses
dependency closure, lifecycle validation, persistence publication, or the
Security Domain policy hook.

This yields one audit rule: the result of an operation depends on the pair
\emph{granted operation, visible target}. Global scope is the natural result
of selecting the unique ancestor of all objects, not an alternative path
around the reference monitor.

%--------------------------------------------------
\subsection{Management Commands versus Application Commands}
\label{sec:gp-management-vs-application}

% Plan:
% - provide a user-visible distinction between GP administration and application use
% - keep it conceptual and operational
% - connect directly with selected context and SCP

For practical use of the secure element, the most important distinction is
between \textbf{management commands} and \textbf{application commands}.

\begin{itemize}
\item \textbf{Management commands} belong to the GlobalPlatform control
  plane. They are used to establish a secure channel, inspect the
  registry, load content, install instances, manage lifecycle states, and
  manipulate key material. They are interpreted under the authority of a
  selected Security Domain.
\item \textbf{Application commands} belong to the business logic of a
  selected application instance. They may use the same APDU transport, but
  their semantics is external to GlobalPlatform secure-element administration.
\end{itemize}

This distinction is operationally essential. A security specialist who
``uses the secure element'' typically alternates between two very
different activities:

\begin{enumerate}
\item \textbf{administering the secure element}, by selecting a Security Domain and
  using authenticated management commands;
\item \textbf{using an application}, by selecting the target application
  AID and sending the application-specific APDUs expected by that
  instance.
\end{enumerate}

The same transport and the same \texttt{SELECT} mechanism are therefore
shared by both planes, but the selected context determines which command
space is actually meaningful.

%--------------------------------------------------
\subsection{Administrative Session Pattern}
\label{sec:gp-admin-session-pattern}

% Plan:
% - give a concrete card-usage mental model to the reader
% - summarize the minimal information needed by an operator
% - avoid implementation details

From an off-card operational perspective, a minimal administrative
session usually requires the following information:

\begin{itemize}
\item the AID of the Security Domain to select;
\item the SCP profile expected by the secure element, such as SCP03 or one of
  the configured SCP11 variants;
\item the relevant key version and static key set associated with that
  Security Domain;
\item the AIDs of the load file and application instance being managed;
\item the privileges and lifecycle expectations attached to the target
  objects.
\end{itemize}

With these elements, the usual administrative pattern is:

\begin{enumerate}
\item select the administrative Security Domain;
\item establish the secure channel;
\item perform the intended management workflow;
\item optionally select the resulting application instance and switch to
  application-level use.
\end{enumerate}

This pattern is simple, but it is central. It is the conceptual bridge
between the APDU and SCP material of Chapter~3 and the more detailed
object, lifecycle, and authority model presented in the present chapter.


%--------------------------------------------------
\section{Consistency and Constraints}
\label{sec:gp-consistency}

% Plan:
% - state the main GP invariants
% - connect them to object relationships and lifecycle model
% - avoid implementation detail while making later design constraints explicit

The GlobalPlatform model and the stricter oXiDe SE tree projection imply the
following implementation requirements:

\begin{itemize}
\item a lifecycle transition must be valid for the addressed object type;
\item a simple deletion cannot leave a live dependent object, while cumulative
  deletion must compute and remove the complete authorized dependency closure;
\item administrative commands must be evaluated under a Security Domain
  that holds the relevant authority;
\item registry updates, deployment steps, and personalization actions
  must preserve a coherent secure-element view of identity, authority, and
  dependency;
\item exactly one oXiDe Root Security Domain has no parent, and it cannot be
  deleted or replaced post-issuance;
\item zero or one Security Domain is explicitly declared as Issuer; when
  present it is a direct child of the root, and a second Issuer declaration is
  rejected;
\item every other managed object has one resolvable Security Domain parent;
\item the parent relation is acyclic, child visibility is a subset of parent
  visibility, and child rights are a subset of parent rights;
\item sibling and detached hierarchies do not reveal objects to one another;
\item a root-resolved DELETE, SET STATUS, or GET STATUS invokes the same
  subtree operation and checks as a subordinate invocation, with the root AID
  as authority;
\item GlobalPlatform global privilege bits do not widen a non-root authority's
  visible subtree in the beta; this is a documented profile restriction.
\end{itemize}

These rules should be understood as conceptual invariants before they are
understood as implementation checks. They explain why some command
sequences are legal and others are not, and they also explain why the
reference implementation later needs explicit registry validation,
privilege checking, and lifecycle guards.

For a reader concerned with operational use, the same invariants can be
read as a checklist of expected failure causes:
\begin{itemize}
\item wrong Security Domain selected;
\item secure channel absent or established under the wrong authority;
\item target object in an incompatible lifecycle state;
\item dependency still present on a load file or application;
\item insufficient privilege for the intended management action.
\end{itemize}

These are not implementation accidents. They are the normal conceptual
reasons why a correctly encoded APDU may still be rejected by the secure
element.

%--------------------------------------------------
\section{Beta Conformance Summary}

The following table separates the oXiDe SE contract from the current beta
claim. ``Implemented'' means that a corresponding code path and regression
test exist; it is not a claim of formal verification.

\begin{description}
\item[Implemented.] The beta has one non-deletable parentless technical root;
  subtree-filtered registry lookup; monotone child privileges; package and
  instance dependency checks; cumulative removal of Security Domain
  descendants and their owned content; explicit zero-or-one Issuer declaration;
  parent-before-child predeployment; self-association projection for direct
  root children; and root-authority reuse of the common DELETE, SET STATUS, and
  GET STATUS operations. Unit and scenario tests cover these internal rules.
\item[Configuration dependent.] A manifest with no \texttt{role = "issuer"}
  has no distinct Issuer Security Domain. Exactly one such declaration creates
  the Issuer below the root; more than one is rejected. Development manifests
  may intentionally omit it, whereas a production-oriented GlobalPlatform
  hierarchy normally includes it.
\item[Partial.] The beta combines installation and selection; it uses a compact
  Application/Security-Domain lifecycle projection; a non-standard package
  lock state; a diagnostic, non-conformant subset of SET STATUS; root-only
  transverse administration rather than the complete GlobalPlatform global
  privilege semantics; and non-atomic multi-entry PUT KEY updates.
\item[Incomplete.] Personalization, extradition, Tokens, DAPs, Receipts, and
  the remaining GlobalPlatform workflow forms are not fully implemented. The
  Card/Issuer lifecycle, standard key deletion/locking management, logically
  deleted load files, and complete cumulative cross-hierarchy reference repair
  are also absent.
\item[Planned post-beta.] Logical channels must provide independent selection
  and Secure Channel context per channel.
\end{description}

The distinction between standard global privileges and oXiDe's current
root-only transverse scope is particularly important for evaluation. The
registry correctly hides its technical containment edge in the GlobalPlatform
association projection, but that projection does not by itself establish full
interoperability with Security Domains that expect Global Delete, Global Lock,
or Global Registry privileges to cross association hierarchies.

%--------------------------------------------------
\section{Discussion}

% Plan:
% - summarize chapter 4 as GP concepts only
% - transition explicitly to chapter 5, which maps these concepts to our system decomposition

This chapter has separated the GlobalPlatform association forest from the
oXiDe SE implementation contract that embeds it below one technical root. The
result is one uniform authorization model: a Security Domain performs granted
operations over its visible subtree; the root is global because its subtree is
the complete registry. The next chapter introduces the software decomposition
that implements this contract, while later implementation and validation
chapters assess which beta claims are supported by code and tests.

%--------------------------------------------------
